% GENERATED FILE. Edit canonical lessons, then run npm run generate:books.
% canonical-source-hash: fnv1a64:b0a9af7d71a6fe94
% canonical-lessons: ML-C189-tuppuka, ML-C189-piliyuka, ML-C189-kulukkuka, ML-C189-uruttuka, ML-C189-ilayuka, ML-R189-first-pass-words-from-chapters-181-184, ML-R189-second-pass-words-from-chapters-185-189

\chapter{To spit, To wring, To shake, To roll, To crawl}
\label{ch:ml-to-spit-to-wring-to-shake-to-roll-to-crawl}
\hlchaptermodality{189}

\begin{chapteropening}
\textbf{By the end of this chapter:} \emph{I can say to spit, to wring, to shake, to roll and to crawl.}

Complete the last lesson of chapter 189: I can say to spit, to wring, to shake, to roll and to crawl.
\end{chapteropening}

\section[\texorpdfstring{\ml{തുപ്പുക} (tuppuka)}{തുപ്പുക (tuppuka)}]{\ml{തുപ്പുക} (tuppuka) — to spit}

\label{lesson:ML-C189-tuppuka}



\begin{quote}
Before the new one: say the Malayalam for to publish, then the Malayalam for to repair.
\end{quote}

\subsection*{You'll want to know: \ml{തുപ്പുക}}
\textbf{\ml{തുപ്പുക}} — \emph{tuppuka} — \textquotedblleft{}to spit\textquotedblright{}.

\begin{grammarlens}[title={What you've built}]
One more word for this chapter.
\end{grammarlens}

\subsection*{Your turn}
\begin{itemize}
\raggedright
  \item \emph{Say it:} \emph{tuppuka}
  \item \emph{Say it:} \emph{tuppuka}, once more
  \item \emph{Say it:} say \emph{tuppuka}
  \item \emph{Recall:} say the Malayalam for to publish, then the Malayalam for to repair, then say \emph{tuppuka} again
\end{itemize}

\begin{tcolorbox}[breakable,colback=teal!4,colframe=teal!35!black,title={Before you move on}]
What does \ml{തുപ്പുക} mean? (\textquotedblleft{}To spit\textquotedblright{}.) Say it once more.
\end{tcolorbox}
\section[\texorpdfstring{\ml{പിഴിയുക} (piḻiyuka)}{പിഴിയുക (piḻiyuka)}]{\ml{പിഴിയുക} (piḻiyuka) — to wring, to squeeze out}

\label{lesson:ML-C189-piliyuka}



\begin{quote}
Before the new one: say the Malayalam for to repair, then the Malayalam for to spit.
\end{quote}

\subsection*{You'll want to know: \ml{പിഴിയുക}}
\textbf{\ml{പിഴിയുക}} — \emph{piḻiyuka} — \textquotedblleft{}to wring, to squeeze out\textquotedblright{}.

\begin{grammarlens}[title={What you've built}]
One more word for this chapter.
\end{grammarlens}

\subsection*{Your turn}
\begin{itemize}
\raggedright
  \item \emph{Say it:} \emph{piḻiyuka}
  \item \emph{Say it:} \emph{piḻiyuka}, once more
  \item \emph{Say it:} say \emph{piḻiyuka}
  \item \emph{Recall:} say the Malayalam for to repair, then the Malayalam for to spit, then say \emph{piḻiyuka} again
\end{itemize}

\begin{tcolorbox}[breakable,colback=teal!4,colframe=teal!35!black,title={Before you move on}]
What does \ml{പിഴിയുക} mean? (\textquotedblleft{}To wring, to squeeze out\textquotedblright{}.) Say it once more.
\end{tcolorbox}
\section[\texorpdfstring{\ml{കുലുക്കുക} (kulukkuka)}{കുലുക്കുക (kulukkuka)}]{\ml{കുലുക്കുക} (kulukkuka) — to shake}

\label{lesson:ML-C189-kulukkuka}



\begin{quote}
Before the new one: say the Malayalam for to spit, then the Malayalam for to wring.
\end{quote}

\subsection*{You'll want to know: \ml{കുലുക്കുക}}
\textbf{\ml{കുലുക്കുക}} — \emph{kulukkuka} — \textquotedblleft{}to shake\textquotedblright{}.

\begin{grammarlens}[title={What you've built}]
One more word for this chapter.
\end{grammarlens}

\subsection*{Your turn}
\begin{itemize}
\raggedright
  \item \emph{Say it:} \emph{kulukkuka}
  \item \emph{Say it:} \emph{kulukkuka}, once more
  \item \emph{Say it:} say \emph{kulukkuka}
  \item \emph{Recall:} say the Malayalam for to spit, then the Malayalam for to wring, then say \emph{kulukkuka} again
\end{itemize}

\begin{tcolorbox}[breakable,colback=teal!4,colframe=teal!35!black,title={Before you move on}]
What does \ml{കുലുക്കുക} mean? (\textquotedblleft{}To shake\textquotedblright{}.) Say it once more.
\end{tcolorbox}
\section[\texorpdfstring{\ml{ഉരുട്ടുക} (uruṭṭuka)}{ഉരുട്ടുക (uruṭṭuka)}]{\ml{ഉരുട്ടുക} (uruṭṭuka) — to roll (something along)}

\label{lesson:ML-C189-uruttuka}



\begin{quote}
Before the new one: say the Malayalam for to wring, then the Malayalam for to shake.
\end{quote}

\subsection*{You'll want to know: \ml{ഉരുട്ടുക}}
\textbf{\ml{ഉരുട്ടുക}} — \emph{uruṭṭuka} — \textquotedblleft{}to roll (something along)\textquotedblright{}.

\begin{grammarlens}[title={What you've built}]
One more word for this chapter.
\end{grammarlens}

\subsection*{Your turn}
\begin{itemize}
\raggedright
  \item \emph{Say it:} \emph{uruṭṭuka}
  \item \emph{Say it:} \emph{uruṭṭuka}, once more
  \item \emph{Say it:} say \emph{uruṭṭuka}
  \item \emph{Recall:} say the Malayalam for to wring, then the Malayalam for to shake, then say \emph{uruṭṭuka} again
\end{itemize}

\begin{tcolorbox}[breakable,colback=teal!4,colframe=teal!35!black,title={Before you move on}]
What does \ml{ഉരുട്ടുക} mean? (\textquotedblleft{}To roll (something along)\textquotedblright{}.) Say it once more.
\end{tcolorbox}
\section[\texorpdfstring{\ml{ഇഴയുക} (iḻayuka)}{ഇഴയുക (iḻayuka)}]{\ml{ഇഴയുക} (iḻayuka) — to crawl}

\label{lesson:ML-C189-ilayuka}



\begin{quote}
Before the new one: say the Malayalam for to shake, then the Malayalam for to roll.
\end{quote}

\subsection*{You'll want to know: \ml{ഇഴയുക}}
\textbf{\ml{ഇഴയുക}} — \emph{iḻayuka} — \textquotedblleft{}to crawl\textquotedblright{}.

\begin{grammarlens}[title={What you've built}]
One more word for this chapter.
\end{grammarlens}

\subsection*{Your turn}
\begin{itemize}
\raggedright
  \item \emph{Say it:} \emph{iḻayuka}
  \item \emph{Say it:} \emph{iḻayuka}, once more
  \item \emph{Say it:} say \emph{iḻayuka}
  \item \emph{Recall:} say the Malayalam for to shake, then the Malayalam for to roll, then say \emph{iḻayuka} again
\end{itemize}

\begin{tcolorbox}[breakable,colback=teal!4,colframe=teal!35!black,title={Before you move on}]
What does \ml{ഇഴയുക} mean? (\textquotedblleft{}To crawl\textquotedblright{}.) Say it once more.
\end{tcolorbox}
\section[(dialogue)]{Words from chapters 181-184 — a first pass}

\label{lesson:ML-R189-first-pass-words-from-chapters-181-184}



\begin{quote}
Say the words for to roll and to crawl.
\end{quote}

\subsection*{Your turn}
Say these aloud:

\begin{itemize}
\raggedright
  \item to lie down, to turn, to cross, to pull, to push
  \item to stir, to fry, to boil, to roast, to pour
  \item to choose, to use, to try, to get, to lose
  \item to break, to get better, to exercise, to fill in, to sign
\end{itemize}

\begin{tcolorbox}[breakable,colback=teal!4,colframe=teal!35!black,title={Before you move on}]
To lie down? (\textbf{\ml{കിടക്കുക}}, \emph{kiṭakkuka}.) To turn? (\textbf{\ml{തിരിയുക}}, \emph{tiriyuka}.) To cross? (\textbf{\ml{കടക്കുക}}, \emph{kaṭakkuka}.) To pull? (\textbf{\ml{വലിക്കുക}}, \emph{valikkuka}.) To push? (\textbf{\ml{തള്ളുക}}, \emph{taḷḷuka}.) To stir? (\textbf{\ml{ഇളക്കുക}}, \emph{iḷakkuka}.) To fry? (\textbf{\ml{പൊരിക്കുക}}, \emph{porikkuka}.) To boil? (\textbf{\ml{തിളപ്പിക്കുക}}, \emph{tiḷappikkuka}.) To roast? (\textbf{\ml{വറുക്കുക}}, \emph{vaṟukkuka}.) To pour? (\textbf{\ml{ഒഴിക്കുക}}, \emph{oḻikkuka}.) To choose? (\textbf{\ml{തിരഞ്ഞെടുക്കുക}}, \emph{tiraññeṭukkuka}.) To use? (\textbf{\ml{ഉപയോഗിക്കുക}}, \emph{upayōgikkuka}.) To try? (\textbf{\ml{ശ്രമിക്കുക}}, \emph{śramikkuka}.) To get? (\textbf{\ml{കിട്ടുക}}, \emph{kiṭṭuka}.) To lose? (\textbf{\ml{നഷ്ടപ്പെടുക}}, \emph{naṣṭappeṭuka}.) To break? (\textbf{\ml{ഒടിയുക}}, \emph{oṭiyuka}.) To get better? (\textbf{\ml{ഭേദമാകുക}}, \emph{bhēdamākuka}.) To exercise? (\textbf{\ml{വ്യായാമം} \ml{ചെയ്യുക}}, \emph{vyāyāmaṁ ceyyuka}.) To fill in? (\textbf{\ml{പൂരിപ്പിക്കുക}}, \emph{pūrippikkuka}.) To sign? (\textbf{\ml{ഒപ്പിടുക}}, \emph{oppiṭuka}.)
\end{tcolorbox}
\section[(dialogue)]{Words from chapters 185-189 — a second pass}

\label{lesson:ML-R189-second-pass-words-from-chapters-185-189}



\begin{quote}
Say the words for to resign and to inform.
\end{quote}

\subsection*{Your turn}
Say these aloud:

\begin{itemize}
\raggedright
  \item to resign, to inform, to charge, to fly, to follow
  \item to share, to hug, to quarrel, to sulk, to tease
  \item to blow, to bloom, to plant, to bark, to rise
  \item to erase, to switch on, to switch off, to publish, to repair
  \item to spit, to wring, to shake, to roll, to crawl
\end{itemize}

\begin{tcolorbox}[breakable,colback=teal!4,colframe=teal!35!black,title={Before you move on}]
To resign? (\textbf{\ml{രാജിവയ്ക്കുക}}, \emph{rājivaykkuka}.) To inform? (\textbf{\ml{അറിയിക്കുക}}, \emph{aṟiyikkuka}.) To charge? (\textbf{\ml{ചാർജ്} \ml{ചെയ്യുക}}, \emph{cārjŭ ceyyuka}.) To fly? (\textbf{\ml{പറക്കുക}}, \emph{paṟakkuka}.) To follow? (\textbf{\ml{പിന്തുടരുക}}, \emph{pintuṭaruka}.) To share? (\textbf{\ml{പങ്കുവെക്കുക}}, \emph{paṅkuvekkuka}.) To hug? (\textbf{\ml{കെട്ടിപ്പിടിക്കുക}}, \emph{keṭṭippiṭikkuka}.) To quarrel? (\textbf{\ml{വഴക്കിടുക}}, \emph{vaḻakkiṭuka}.) To sulk? (\textbf{\ml{പിണങ്ങുക}}, \emph{piṇaṅṅuka}.) To tease? (\textbf{\ml{കളിയാക്കുക}}, \emph{kaḷiyākkuka}.) To blow? (\textbf{\ml{വീശുക}}, \emph{vīśuka}.) To bloom? (\textbf{\ml{വിരിയുക}}, \emph{viriyuka}.) To plant? (\textbf{\ml{നടുക}}, \emph{naṭuka}.) To bark? (\textbf{\ml{കുരയ്ക്കുക}}, \emph{kuraykkuka}.) To rise? (\textbf{\ml{ഉദിക്കുക}}, \emph{udikkuka}.) To erase? (\textbf{\ml{മായ്ക്കുക}}, \emph{māykkuka}.) To switch on? (\textbf{\ml{ഓണാക്കുക}}, \emph{ōṇākkuka}.) To switch off? (\textbf{\ml{ഓഫാക്കുക}}, \emph{ōphākkuka}.) To publish? (\textbf{\ml{പ്രസിദ്ധീകരിക്കുക}}, \emph{prasiddhīkarikkuka}.) To repair? (\textbf{\ml{നന്നാക്കുക}}, \emph{nannākkuka}.) To spit? (\textbf{\ml{തുപ്പുക}}, \emph{tuppuka}.) To wring? (\textbf{\ml{പിഴിയുക}}, \emph{piḻiyuka}.) To shake? (\textbf{\ml{കുലുക്കുക}}, \emph{kulukkuka}.) To roll? (\textbf{\ml{ഉരുട്ടുക}}, \emph{uruṭṭuka}.) To crawl? (\textbf{\ml{ഇഴയുക}}, \emph{iḻayuka}.)
\end{tcolorbox}
